\documentclass[14pt,letterpaper]{extarticle}
\usepackage[utf8]{inputenc}

% Set margins
% https://tex.stackexchange.com/a/327136
% this should replace use of the 'fullpage' package
\usepackage[
    includehead,
    margin=2.5cm,
    bmargin=2cm,
]{geometry}

\usepackage{amssymb}
\usepackage{amsthm}
\usepackage{verbatim} % This works better than comment for some reason
\usepackage{hyperref}
\usepackage{amsmath}
\usepackage{gensymb}
\usepackage{enumitem}
\usepackage{xfrac} % Fractions - https://tex.stackexchange.com/a/230973
\usepackage{xcolor}

% The magic latex drawing thing
\usepackage{tikz}
\usetikzlibrary{snakes,arrows,shapes}
\usepackage[en-US,showdow]{datetime2}

% Colored hyperlinks:
% https://www.overleaf.com/learn/latex/Hyperlinks
\hypersetup{
    colorlinks=true,
    urlcolor=blue
}

% Setup nice headings
% https://www.overleaf.com/learn/latex/Headers_and_footers
\usepackage{fancyhdr}
\usepackage{titling} % Save title for use in fancyhdr

% Shortcuts
\newcommand*{\bP}{\mathbf{P}}
\newcommand*{\bq}{\mathbf{q}}


% Use black tombstone for QED
\renewcommand{\qedsymbol}{$\blacksquare$}

% Important sets
\newcommand*{\C}{\mathbb{C}}
\newcommand*{\R}{\mathbb{R}}
\newcommand*{\Q}{\mathbb{Q}}
\newcommand*{\Z}{\mathbb{Z}}

%
% symbols
%

% name by meaning, not by character
\newcommand{\grad}{\nabla}


% New command for 'partial QED' using white tombstone
\newcommand*{\partialqed}{$\hfill \square$ \\}

\newcommand*{\todo}{{\color{red}\textbf{TODO:}} }

% Increased spaces spacing
\usepackage{setspace} \onehalfspacing

\title{Math 468 - Homework 6}
\author{Nicholas Schlabach (nickninja@arizona.edu)}
\DTMsavedate{titledate}{2026-03-24}
\date{\DTMUsedate{titledate}}

%% Copied from https://tex.stackexchange.com/a/2238
%% Modified by Techcable 9/19/21 to use '[...]' instead of '(...)'
\newenvironment{amatrix}[1]{%
  \left[\begin{array}{@{}*{#1}{c}|c@{}}
}{%
  \end{array}\right]
}

\begin{document}
\pagestyle{fancy}
\fancyhead[L]{Nicholas Schlabach}
\fancyhead[C]{\hspace{5mm}\Big(\thetitle{}\Big)}
\fancyhead[R]{{\DTMsetdatestyle{mmddyy}\DTMusedate{titledate}}, Page \thepage}
\fancyfoot{} % no footer
\maketitle
% maketitle automatically sets pagestyle=plain, override this
\thispagestyle{fancy}


\begin{center}
\begin{enumerate}
\setcounter{enumi}{1}
\item
    \begin{enumerate}
    \item
        Observe that $\langle1,1\rangle$ is an eigenvector with eigenvalue $\lambda=0$.
        Solving for $\det( \bq-\lambda I)=0$ gives characteristic polynomial $\lambda(\lambda_0+\mu_1+\lambda)=0$.
        The eigenvector $\langle \lambda_0,-\mu_1\rangle$ is associated with eigenvalue $\lambda=-\lambda_0-\mu_1$.
    \item
        Diagonalizing using the eigenvectors gives 
        \[ Q=\begin{bmatrix}
            1 & \lambda_0 \\
            1 & -\mu_1
        \end{bmatrix}, \; Q^{-1}\bq
            Q=\begin{bmatrix}
            0 & 0 \\
            0 & -\lambda_0-\mu_1
        \end{bmatrix}=\boldsymbol{\lambda}, Q^{-1}=\frac{1}{\lambda+\mu}\begin{bmatrix}
                \mu & \lambda \\
                1 & -1
        \end{bmatrix} \]
        \item The Wikipedia page for matrix exponential says that if $Y$ is invertible, then $e^{YXY{-1}}=Ye^XY^{-1}$. Along with the fact that $t\bq$ is diagonalizable whenever $\bq$ is we get
        \[ \bP(t)=e^{t\bq}=Qe^{t\boldsymbol{\lambda}}Q^{-1}=Q\begin{bmatrix}
            1 & 0 \\
            0 & e^{t(-\lambda_0-\mu_1)}       \end{bmatrix}Q^{-1} \]
        Multiplying these matrices gives the final result
        \[ \bP(t)=e^{t\bq}=\frac{1}{\lambda_0+\mu_1}\begin{bmatrix}
                \lambda_0e^{-t(\lambda_0+\mu_1)}+\mu_1 & \lambda_0\big(1-e^{-t(\lambda_0+\mu_1)}\big) \\
                \mu_1\big(1-e^{-t(\lambda_0+\mu_1)}\big) & \mu_1 e^{-t(\lambda_0+\mu_1)}+\lambda_0
            \end{bmatrix} \]
        \item
            Using the fact that $\lambda_0,\mu_1>0$ and $\forall k>0,\lim_{t \to \infty}e^{-tk}=0$ we get
            \[ \lim_{t \to \infty} \bP(t)=\frac{1}{\lambda_0+\mu_1}\begin{bmatrix}
                \mu_1 & \lambda_0 \\
                \mu_1 & \lambda_0
            \end{bmatrix}\]
    \end{enumerate}
\end{enumerate}
\end{center}
\end{document}